\documentclass[preprint]{article}
\usepackage{neurips_2025}
\usepackage[utf8]{inputenc}
\usepackage[T1]{fontenc}
\usepackage{hyperref}
\usepackage{url}
\usepackage{booktabs}
\usepackage{amsfonts}
\usepackage{amsmath}
\usepackage{nicefrac}
\usepackage{microtype}
\usepackage{graphicx}
\usepackage{xcolor}
\usepackage{algorithm}
\usepackage{algorithmic}
\usepackage{multirow}

\title{MarkovTier: Anticipatory Page Migration via\\Markov Phase Models for Tiered Memory Systems}

\author{
  Anonymous Author(s)
}

\begin{document}

\maketitle

%% ============================================================
%% ABSTRACT
%% ============================================================
\begin{abstract}
Modern tiered memory systems combining fast DRAM with slower CXL-attached memory rely on page migration policies to place hot pages in fast memory. All existing policies---including TPP, Nomad, ALTO, and ARMS---are fundamentally \emph{reactive}: they detect access pattern changes after they occur and then migrate, incurring a \emph{migration lag} during which applications suffer degraded performance. We observe that program execution exhibits predictable phase behavior and propose \textbf{MarkovTier}, a framework that models workload phase transitions as a first-order Markov chain and \emph{proactively} migrates pages before phase transitions occur. MarkovTier combines online phase detection via working set signatures with a learned transition probability model that triggers anticipatory migration when prediction confidence is high. In trace-driven simulations across seven workload types with ten random seeds, MarkovTier reduces the overall slow-tier access ratio by up to 40\% on workloads with predictable phase structure compared to the best reactive baseline, while performing comparably on workloads with irregular phases. MarkovTier also reduces migration lag by 26--39\% and increases total migration volume by only 11\%. We show that the improvement stems primarily from Markov-guided anticipatory migration rather than phase-aware reactive migration alone, and discuss the limitations of first-order models on workloads with complex transition patterns.
\end{abstract}

%% ============================================================
%% 1. INTRODUCTION
%% ============================================================
\section{Introduction}

Tiered memory architectures that combine fast local DRAM with higher-capacity but slower CXL-attached memory are increasingly deployed in modern data centers~\citep{maruf2023tpp, xiang2024nomad, liu2025hotness}. The performance of these systems depends critically on \emph{page migration policies} that place frequently accessed (``hot'') pages in fast memory while relegating infrequently accessed (``cold'') pages to slower tiers.

All existing migration policies are fundamentally \textbf{reactive}: they observe page access patterns, detect hot or cold pages \emph{after} a significant number of accesses have occurred, and then initiate migrations. This reactive approach incurs an inherent \textbf{migration lag}---the time interval between when a workload's access pattern changes (e.g., due to a program phase transition) and when affected pages have been migrated to the appropriate tier. During this lag, newly-hot pages reside in slow memory while newly-cold pages wastefully occupy fast memory, degrading application performance.

This migration lag is not merely theoretical. MTM~\citep{xu2024mtm} demonstrated that DAMON-based profiling takes 3$\times$ longer than their approach to achieve 50\% profiling accuracy after a phase change. ARMS~\citep{yadalam2025arms} introduced change-point detection to accelerate reaction to phase shifts but still fundamentally reacts to changes after they manifest. Even the state-of-the-art ALTO system~\citep{liu2025hotness} remains reactive in its core design.

\textbf{Key insight.} Programs exhibit well-documented \emph{phase behavior}: they cycle through distinct execution phases with different memory access patterns~\citep{sherwood2002simpoint, dhodapkar2002working}. Crucially, phase transitions in many workloads are \emph{predictable}---programs cycle through phases in repeatable sequences. Yet no tiered memory system exploits this predictability for proactive migration.

We propose \textbf{MarkovTier}, a phase-aware anticipatory page migration framework that models workload phase transitions as a Markov chain and proactively migrates pages \emph{before} phase transitions occur. Our contributions are:
\begin{itemize}
    \item We propose \textbf{Markov-guided anticipatory migration}, the first use of predictive phase transition models to drive proactive page migration in tiered memory systems, transforming the paradigm from reactive detect-then-migrate to predictive anticipate-then-migrate.
    \item We introduce \textbf{phase-aware reactive migration}, where detected phase profiles accelerate reactive migration during and after transitions, providing benefits even when anticipatory prediction is not possible.
    \item We conduct a comprehensive evaluation across seven workload types and three adversarial scenarios with ten random seeds each, demonstrating up to 40\% reduction in slow-tier access ratio on predictable workloads with no degradation on unpredictable ones, and analyze the conditions under which anticipatory migration succeeds and fails.
\end{itemize}

The remainder of this paper is organized as follows. Section~\ref{sec:related} reviews related work. Section~\ref{sec:method} describes MarkovTier's design. Section~\ref{sec:experiments} presents our experimental evaluation. Section~\ref{sec:discussion} discusses limitations. Section~\ref{sec:conclusion} concludes.

%% ============================================================
%% 2. RELATED WORK
%% ============================================================
\section{Related Work}
\label{sec:related}

\paragraph{Reactive tiered memory management.}
TPP~\citep{maruf2023tpp} introduced transparent page placement for CXL-enabled tiered memory in the Linux kernel, using NUMA hinting faults for promotion and page reclamation for demotion. Nomad~\citep{xiang2024nomad} proposed non-exclusive memory tiering via transactional page migration, maintaining shadow copies to reduce demotion cost. ALTO~\citep{liu2025hotness} moved beyond simple hotness to use amortized offcore latency (AOL) as a migration metric, throttling unnecessary migrations. MTM~\citep{xu2024mtm} improved profiling speed over DAMON but still detects phases after they occur. Nimble~\citep{yan2019nimble} optimized the page migration mechanism itself to reduce overhead. All these systems share a common limitation: they react to access pattern changes rather than anticipating them.

\paragraph{Adaptive tiered memory.}
ARMS~\citep{yadalam2025arms} introduced change-point detection on slow-tier bandwidth to detect hot set shifts, adapting profiling window length based on phase stability. ARMS is the closest existing work to ours in spirit, as it adapts behavior based on phase awareness. However, ARMS reacts to detected changes rather than anticipating future phases. FlexMem~\citep{xu2024flexmem} combines different profiling methods to respond to application phase changes using adaptive profiling granularity. ArtMem~\citep{yi2025artmem} applies Q-learning to learn adaptive migration policies. While ArtMem's RL agent implicitly captures some phase behavior, it does not explicitly model phase transitions or perform anticipatory migration.

\paragraph{Program phase detection.}
SimPoint~\citep{sherwood2002simpoint} demonstrated that programs exhibit phase behavior identifiable through Basic Block Vectors, establishing that phase structure is pervasive and predictable. \citet{dhodapkar2002working} proposed working set signatures for dynamic phase detection. Our approach extends these techniques to the memory page level for tiering decisions, adding a Markov transition model for prediction---a step not taken by prior phase detection work in the memory management context.

\paragraph{Positioning.}
Unlike all prior tiered memory systems, which react to phase changes after they occur, MarkovTier explicitly models the sequence of phase transitions as a Markov chain and uses this model to proactively migrate pages before they are needed. This transforms migration timing from a reactive paradigm to a predictive one.

%% ============================================================
%% 3. METHOD
%% ============================================================
\section{Method}
\label{sec:method}

\subsection{Overview}

MarkovTier augments a reactive migration backbone with two novel components: (1) an online phase detector that identifies execution phases via working set signatures and (2) a Markov transition model that predicts imminent phase transitions and triggers anticipatory migration. Figure~\ref{fig:motivation} illustrates the benefit: during phase transitions, MarkovTier's anticipatory migrations reduce slow-tier accesses compared to reactive baselines.

\subsection{Online Phase Detection}
\label{sec:phase_detection}

At each profiling interval (every $I$ memory accesses), MarkovTier monitors the set of accessed pages. Phase transitions are detected using a \emph{containment-based} approach: we compute the fraction of the current short window's pages that appear in the preceding long window. Let $S_\text{short}$ be the set of pages accessed in the last $w_s = 2$ intervals and $S_\text{long}$ the pages in the preceding $w_l = 8$ intervals. We detect a transition when:
\begin{equation}
    \text{containment}(S_\text{short}, S_\text{long}) = \frac{|S_\text{short} \cap S_\text{long}|}{|S_\text{short}|} < \tau_t
\end{equation}
where $\tau_t = 0.55$ is the transition threshold. This containment metric is more sensitive than Jaccard similarity for detecting working set shifts, as it specifically measures whether the new pages have departed from the previous set.

Phase identification uses top-$K$ page Jaccard similarity against stored phase signatures. When a transition is detected, the current signature is compared to all known phases; if the best match exceeds a merge threshold ($\tau_m = 0.25$), the interval is assigned to that phase and the stored profile is updated via exponential moving average. Otherwise, a new phase is created. This allows MarkovTier to recognize recurring phases across the workload's execution.

\subsection{Markov Transition Model}
\label{sec:markov}

Phase transitions are modeled as a first-order Markov chain over identified phases $\{P_1, \ldots, P_k\}$. The transition matrix $T$ is maintained online:
\begin{equation}
    T[i][j] = \frac{\text{count}(P_i \to P_j)}{\sum_{j'} \text{count}(P_i \to P_{j'})}
\end{equation}
For each phase $P_i$, we maintain a \emph{page hotness profile}---the per-page access frequency, updated via exponential moving average (EMA, $\alpha = 0.15$) at each interval the phase is active. Hot pages are defined as the top 30\% by EMA-smoothed access count.

Given the current phase $P_i$, the predicted next phase is $P_{j^*} = \arg\max_j T[i][j]$ with confidence $c = T[i][j^*]$.

\subsection{Anticipatory Migration}
\label{sec:anticipatory}

MarkovTier triggers anticipatory migration when two conditions are met: (1) at least one phase transition has been observed (so transition probabilities exist), and (2) the current phase has progressed past a fraction $\phi_s = 0.4$ of its estimated duration. The estimated duration uses the historical mean for the current phase.

As the phase progresses beyond $\phi_s$, the anticipatory migration budget ramps up:
\begin{equation}
    B_\text{antic} = \left\lfloor B_\text{total} \cdot f_\text{antic} \cdot \min\!\left(1, \frac{\phi - \phi_s}{1 - \phi_s}\right) \cdot c \right\rfloor
\end{equation}
where $B_\text{total}$ is the total migration bandwidth per interval, $f_\text{antic} = 0.4$ is the maximum anticipatory fraction, $\phi$ is the current phase progress, and $c$ is the prediction confidence. Pages are selected for anticipatory promotion from the predicted phase's hot page ranking, prioritized by expected access frequency. Speculatively promoted pages are \emph{protected} from reactive demotion until the next transition.

\subsection{Reactive Backbone and Phase-Aware Migration}

The remaining bandwidth ($B_\text{total} - B_\text{antic}$) is used for reactive migration: promoting the most-accessed slow-tier pages and demoting LRU fast-tier pages. This ensures MarkovTier performs no worse than a reactive policy even when predictions are wrong.

Additionally, when a phase transition is detected, MarkovTier immediately prioritizes migration of pages known to be hot in the detected phase (from stored profiles), rather than waiting to re-discover them through access monitoring. We evaluate this \emph{phase-aware reactive} component in isolation in Section~\ref{sec:ablation}.

%% ============================================================
%% 4. EXPERIMENTS
%% ============================================================
\section{Experiments}
\label{sec:experiments}

\subsection{Experimental Setup}

\paragraph{Simulator.} We build a trace-driven tiered memory simulator modeling a two-tier system: fast tier (DRAM, 100\,ns latency) and slow tier (CXL memory, 350\,ns latency). Migration bandwidth is 50 page migrations per profiling interval (1{,}000 accesses). Fast-tier capacity defaults to 50\% of unique pages.

\paragraph{Workloads.} We evaluate on seven workload types spanning a spectrum of phase predictability:
\begin{itemize}
    \item \textbf{Regular}: 5 phases, deterministic cyclic transitions, 50K accesses/phase, 20\% overlap, 1000 pages/phase.
    \item \textbf{Semi-Regular}: Same as Regular but stochastic transitions (80\% dominant cycle, 20\% random).
    \item \textbf{Irregular}: 8 phases, variable duration (20K--80K), uniform random transitions, 30\% overlap.
    \item \textbf{GCC-like}: 10 phases, log-normal durations, 70\% dominant transitions, 25\% overlap.
    \item \textbf{MCF-like}: 4 phases, cyclic, large working sets (3000 pages/phase), 10\% overlap.
    \item \textbf{Webserver}: 3 diurnal phases, highly regular, 60\% overlap, 2000 pages/phase.
    \item \textbf{MapReduce}: 2-phase alternation, 5\% overlap, variable duration.
\end{itemize}
We also test three adversarial scenarios: \textbf{Chaotic} (uniform random, 8 phases), \textbf{Drifting} (gradual change, 85\% overlap), and \textbf{Bursty} (2K--200K variable bursts).

\paragraph{Baselines.}
\textbf{LRU-Reactive}: promotes most-accessed slow pages, demotes LRU (models NUMA balancing).
\textbf{TPP-like}: NUMA fault promotion, periodic demotion scan~\citep{maruf2023tpp}.
\textbf{ARMS-like}: change-point detection with adaptive windows~\citep{yadalam2025arms}.
\textbf{ALTO-like}: AOL-weighted migration with throttling~\citep{liu2025hotness}; our strongest reactive baseline.
\textbf{Oracle}: offline optimal with perfect future knowledge.

\paragraph{Statistical methodology.} All experiments use 10 random seeds. We report mean $\pm$ standard deviation and use Welch's $t$-test for significance testing.

\subsection{Main Results}
\label{sec:main_results}

\begin{table}[t]
\caption{Overall slow-tier access ratio ($\downarrow$ lower is better). Mean $\pm$ std over 10 seeds. Best non-oracle in \textbf{bold}. $^{***}$\!: $p < 0.001$ vs.\ ALTO-like.}
\label{tab:main}
\centering
\footnotesize
\setlength{\tabcolsep}{4pt}
\begin{tabular}{@{}lccccccc@{}}
\toprule
 & Regular & Semi-Reg. & Irregular & GCC & MCF & Web & MR \\
\midrule
LRU-React. & .067{\scriptsize$\pm$.001} & .058{\scriptsize$\pm$.009} & .021{\scriptsize$\pm$.003} & .029{\scriptsize$\pm$.008} & .081{\scriptsize$\pm$.001} & .113{\scriptsize$\pm$.001} & .112{\scriptsize$\pm$.007} \\
TPP-like & .114{\scriptsize$\pm$.002} & .095{\scriptsize$\pm$.014} & .047{\scriptsize$\pm$.002} & .049{\scriptsize$\pm$.006} & .129{\scriptsize$\pm$.002} & .165{\scriptsize$\pm$.002} & .115{\scriptsize$\pm$.003} \\
ARMS-like & .075{\scriptsize$\pm$.001} & .071{\scriptsize$\pm$.005} & .028{\scriptsize$\pm$.002} & .042{\scriptsize$\pm$.005} & .108{\scriptsize$\pm$.001} & .180{\scriptsize$\pm$.002} & .113{\scriptsize$\pm$.004} \\
ALTO-like & .065{\scriptsize$\pm$.001} & .056{\scriptsize$\pm$.009} & .021{\scriptsize$\pm$.003} & .028{\scriptsize$\pm$.008} & .078{\scriptsize$\pm$.001} & .110{\scriptsize$\pm$.001} & .108{\scriptsize$\pm$.007} \\
\textbf{MarkovTier} & \textbf{.039}{\scriptsize$\pm$.002}$^{***}$ & \textbf{.040}{\scriptsize$\pm$.004}$^{***}$ & .022{\scriptsize$\pm$.004} & .027{\scriptsize$\pm$.006} & .078{\scriptsize$\pm$.003} & .113{\scriptsize$\pm$.001} & .099{\scriptsize$\pm$.009} \\
Oracle & .019{\scriptsize$\pm$.001} & .017{\scriptsize$\pm$.002} & .002{\scriptsize$\pm$.000} & .003{\scriptsize$\pm$.000} & .063{\scriptsize$\pm$.001} & .101{\scriptsize$\pm$.001} & .054{\scriptsize$\pm$.002} \\
\bottomrule
\end{tabular}
\end{table}

Table~\ref{tab:main} presents the main results. MarkovTier achieves statistically significant improvements over ALTO-like on the two most predictable workloads: \textbf{40\% reduction on Regular} ($p < 0.001$) and \textbf{29\% reduction on Semi-Regular} ($p < 0.001$). On the remaining workloads---where phase transitions are inherently less predictable---MarkovTier performs comparably to ALTO-like, with no statistically significant degradation on any workload except a marginal 2.5\% increase on Webserver ($p < 0.001$), which we attribute to the high inter-phase overlap (60\%) making anticipatory migrations less targeted.

\begin{figure}[t]
\centering
\includegraphics[width=0.95\linewidth]{figures/fig2_main_comparison.pdf}
\caption{Overall slow-tier access ratio across all policies and workloads (10 seeds). MarkovTier provides the largest improvement on Regular and Semi-Regular workloads.}
\label{fig:main_comparison}
\end{figure}

Figure~\ref{fig:main_comparison} visualizes these results. The benefit is most pronounced on workloads where the Markov model can learn reliable transition probabilities from the deterministic or dominant-path cyclic structure.

\paragraph{Transition-window analysis.}

\begin{table}[t]
\caption{Mean slow-tier ratio during phase transitions (10-interval window). $\downarrow$ lower is better.}
\label{tab:transition}
\centering
\footnotesize
\setlength{\tabcolsep}{4pt}
\begin{tabular}{@{}lccccccc@{}}
\toprule
 & Regular & Semi-Reg. & Irregular & GCC & MCF & Web & MR \\
\midrule
ALTO-like & .239{\scriptsize$\pm$.004} & .205{\scriptsize$\pm$.033} & .105{\scriptsize$\pm$.016} & .122{\scriptsize$\pm$.040} & .289{\scriptsize$\pm$.004} & .295{\scriptsize$\pm$.005} & .309{\scriptsize$\pm$.016} \\
MarkovTier & \textbf{.142}{\scriptsize$\pm$.008} & \textbf{.144}{\scriptsize$\pm$.015} & .109{\scriptsize$\pm$.018} & .117{\scriptsize$\pm$.029} & .275{\scriptsize$\pm$.016} & .301{\scriptsize$\pm$.006} & .259{\scriptsize$\pm$.014} \\
\bottomrule
\end{tabular}
\end{table}

Table~\ref{tab:transition} focuses on the transition windows where anticipatory migration has the most impact. MarkovTier reduces the transition slow-tier ratio by \textbf{41\% on Regular} and \textbf{30\% on Semi-Regular} compared to ALTO-like.

\paragraph{Migration lag and overhead.}

\begin{figure}[t]
\centering
\includegraphics[width=0.95\linewidth]{figures/fig3_migration_lag.pdf}
\caption{Migration lag (intervals to reach $<$10\% slow-tier ratio after transition). MarkovTier reduces lag by 39\% on Regular and 26\% on Semi-Regular workloads.}
\label{fig:migration_lag}
\end{figure}

Figure~\ref{fig:migration_lag} shows migration lag. MarkovTier reduces lag from 9.3 to 5.7 intervals on Regular (39\% reduction) and from 8.1 to 6.0 on Semi-Regular (26\%). On other workloads, lag is comparable to reactive baselines.

\begin{table}[t]
\caption{Total page migrations per workload. Higher values indicate more migration activity.}
\label{tab:migrations}
\centering
\footnotesize
\setlength{\tabcolsep}{4pt}
\begin{tabular}{@{}lccccccc@{}}
\toprule
 & Regular & Semi-Reg. & Irregular & GCC & MCF & Web & MR \\
\midrule
LRU-React. & 20{,}629 & 17{,}867 & 5{,}603 & 8{,}704 & 32{,}565 & 38{,}524 & 24{,}987 \\
ALTO-like & 20{,}627 & 17{,}865 & 5{,}603 & 8{,}707 & 32{,}519 & 38{,}522 & 24{,}975 \\
ARMS-like & 34{,}806 & 33{,}922 & 11{,}846 & 19{,}526 & 39{,}460 & 38{,}529 & 29{,}394 \\
MarkovTier & 22{,}920 & 20{,}146 & 8{,}187 & 10{,}714 & 33{,}363 & 38{,}525 & 27{,}761 \\
\bottomrule
\end{tabular}
\end{table}

Table~\ref{tab:migrations} shows migration overhead. MarkovTier increases migrations by 11\% on Regular compared to ALTO-like---substantially less than ARMS-like's 69\% increase. The additional migrations represent the cost of anticipatory speculation, which is modest relative to the 40\% performance improvement.

\begin{figure}[t]
\centering
\includegraphics[width=0.95\linewidth]{figures/fig1_motivation.pdf}
\caption{Slow-tier access ratio over time for Regular workload (seed 42). Vertical lines mark phase transitions. MarkovTier recovers faster after each transition due to anticipatory migration.}
\label{fig:motivation}
\end{figure}

\paragraph{Adversarial robustness.}
On adversarial workloads, MarkovTier matches ALTO-like: Chaotic (0.036$\pm$0.007 vs.\ 0.036$\pm$0.007), Drifting (0.005$\pm$0.000 vs.\ 0.005$\pm$0.000), Bursty (0.023$\pm$0.006 vs.\ 0.023$\pm$0.004). When transitions are unpredictable, the Markov model's confidence is low and anticipatory migration is suppressed, causing graceful degradation to the reactive backbone.

\subsection{Ablation Study}
\label{sec:ablation}

We evaluate four ablation variants to isolate component contributions:
\begin{itemize}
    \item \textbf{Phase-Only}: Phase detection with stored profiles for fast reactive migration, but no prediction or anticipation.
    \item \textbf{No-Confidence}: Full MarkovTier with confidence threshold = 0 (always anticipate).
    \item \textbf{No-Rollback}: Full MarkovTier with rollback disabled.
    \item \textbf{Random-Anticipation}: Random phase prediction instead of Markov model.
\end{itemize}

\begin{table}[t]
\caption{Ablation study: overall slow-tier access ratio ($\downarrow$). Mean $\pm$ std over 10 seeds.}
\label{tab:ablation}
\centering
\small
\begin{tabular}{lcccc}
\toprule
 & Regular & Semi-Regular & Irregular & GCC-like \\
\midrule
ALTO-like (reference) & .065$\pm$.001 & .056$\pm$.009 & .021$\pm$.003 & .028$\pm$.008 \\
Phase-Only & .066$\pm$.001 & .056$\pm$.009 & .021$\pm$.003 & .029$\pm$.008 \\
Random-Anticipation & .061$\pm$.001 & .052$\pm$.008 & .020$\pm$.003 & .028$\pm$.008 \\
No-Confidence & .039$\pm$.002 & .040$\pm$.004 & .022$\pm$.004 & .027$\pm$.006 \\
No-Rollback & .039$\pm$.002 & .040$\pm$.004 & .022$\pm$.004 & .027$\pm$.006 \\
\textbf{MarkovTier (full)} & \textbf{.039$\pm$.002} & \textbf{.040$\pm$.004} & .022$\pm$.004 & .027$\pm$.006 \\
\bottomrule
\end{tabular}
\end{table}

\begin{figure}[t]
\centering
\includegraphics[width=0.95\linewidth]{figures/fig4_ablation.pdf}
\caption{Ablation study on four representative workloads.}
\label{fig:ablation}
\end{figure}

Table~\ref{tab:ablation} and Figure~\ref{fig:ablation} reveal the component contributions:

\textbf{Markov prediction is the primary contributor.} Phase-Only (0.066 on Regular) performs no better than ALTO-like (0.065), while full MarkovTier (0.039) provides a 40\% reduction. This demonstrates that anticipatory migration---not phase-aware reactive migration---drives the improvement.

\textbf{Structured prediction matters.} Random-Anticipation (0.061) performs much worse than MarkovTier (0.039) on Regular, confirming that the Markov model's learned transition probabilities are essential. Random speculation moves some useful pages by chance but wastes most of its budget.

\textbf{Confidence thresholding and rollback are inactive.} No-Confidence and No-Rollback produce identical results to the full system. The confidence threshold ($\geq 0.3$) is almost always met when predictions are available, and mispredictions are rare enough that rollback provides no measurable benefit. We report this transparently: while these are principled design choices, they are not active contributors under our evaluated conditions. The system's actual effectiveness rests on two innovations: \emph{Markov-guided anticipatory migration} (the dominant contributor) and the \emph{reactive backbone with phase-aware profiles} (the safety net).

\subsection{Sensitivity Analysis}

\begin{figure}[t]
\centering
\includegraphics[width=0.7\linewidth]{figures/fig5_capacity_sensitivity.pdf}
\caption{Slow-tier ratio vs.\ fast-tier capacity on Regular workload. MarkovTier's advantage is largest at 30--60\% capacity.}
\label{fig:capacity}
\end{figure}

\paragraph{Fast-tier capacity.}
Figure~\ref{fig:capacity} shows MarkovTier's benefit across capacity ratios from 10\% to 80\%. The improvement is most pronounced at moderate ratios (30--60\%), where intelligent placement matters most. At very low capacity, the working set far exceeds fast tier and anticipation cannot help; at high capacity, most pages already fit regardless of policy.

\begin{figure}[t]
\centering
\begin{minipage}{0.48\linewidth}
\includegraphics[width=\linewidth]{figures/fig6_confidence_sweep.pdf}
\end{minipage}
\hfill
\begin{minipage}{0.48\linewidth}
\includegraphics[width=\linewidth]{figures/fig7_budget_sweep.pdf}
\end{minipage}
\caption{Sensitivity to confidence threshold (left) and anticipatory budget fraction (right). Both show broad insensitivity across a wide range.}
\label{fig:sensitivity}
\end{figure}

\paragraph{Hyperparameters.}
Figure~\ref{fig:sensitivity} shows that MarkovTier is broadly insensitive to confidence threshold and budget fraction. The confidence threshold has minimal effect because transition probabilities are typically either very high or very low. The budget fraction shows a mild optimum around 0.3--0.5.

\subsection{Second-Order Markov Model}
\label{sec:second_order}

We implemented a second-order Markov model conditioning on the last two phases, $P(P_{t+1} \mid P_t, P_{t-1})$, to test whether higher-order models improve prediction on complex workloads.

\begin{table}[t]
\caption{First-order vs.\ second-order Markov: overall slow-tier ratio ($\downarrow$). Mean $\pm$ std over 10 seeds.}
\label{tab:second_order}
\centering
\small
\begin{tabular}{lcccc}
\toprule
 & Regular & Semi-Regular & Irregular & GCC-like \\
\midrule
1st order & .039$\pm$.002 & .040$\pm$.004 & .022$\pm$.004 & .027$\pm$.006 \\
2nd order & .039$\pm$.002 & .040$\pm$.004 & .022$\pm$.004 & .027$\pm$.006 \\
\bottomrule
\end{tabular}
\end{table}

Table~\ref{tab:second_order} shows that the second-order model provides minimal improvement. On Regular workloads, first-order is already optimal since transitions depend only on the current phase. On GCC-like workloads with 10 phases, the bigram state space ($10 \times 10 = 100$) far exceeds the number of observed transitions ($\sim$12), leading to sparse estimates that are no more accurate than first-order. The fundamental challenge on complex workloads is inherent unpredictability, not model order.

%% ============================================================
%% 5. DISCUSSION
%% ============================================================
\section{Discussion and Limitations}
\label{sec:discussion}

\paragraph{When does MarkovTier help?}
MarkovTier provides substantial benefits when workloads have predictable phase transitions and moderate fast-tier capacity (30--60\%). On workloads with random transitions (Irregular, Chaotic) or high working set overlap (Webserver), MarkovTier performs comparably to reactive baselines---the system gracefully degrades to its reactive backbone.

\paragraph{Inactive components.}
Confidence thresholding and rollback are principled but empirically inactive. The confidence threshold is rarely binding because transition probabilities are typically either very high or uniformly low. Rollback is inactive because mispredictions are rare when anticipation is triggered. More diverse workloads with intermediate predictability may activate these mechanisms.

\paragraph{Simulation limitations.}
Our evaluation uses trace-driven simulation, which abstracts away real-system effects including OS scheduling, TLB shootdowns, NUMA topology, cache effects, and memory contention. Real-system validation on CXL hardware is an important future direction.

\paragraph{First-order limitations.}
The first-order Markov assumption limits effectiveness on workloads with history-dependent transitions. Our second-order experiment showed that simply increasing model order does not help due to data sparsity. Variable-order models or neural predictors may help at increased complexity.

\paragraph{Migration overhead.}
MarkovTier increases migration volume by 11\% on Regular workloads. In bandwidth-constrained deployments, this overhead could interfere with application traffic. The budget fraction parameter provides a knob to control this tradeoff.

%% ============================================================
%% 6. CONCLUSION
%% ============================================================
\section{Conclusion}
\label{sec:conclusion}

We proposed MarkovTier, the first tiered memory migration framework to exploit predictable phase transitions through Markov-guided anticipatory page migration. By proactively migrating pages before transitions occur, MarkovTier reduces slow-tier accesses by up to 40\% on workloads with predictable phase structure, while maintaining parity on unpredictable workloads. Migration lag is reduced by 26--39\% at a cost of 11\% additional migration volume.

Our ablation study shows the improvement stems from anticipatory migration guided by learned Markov transition probabilities, not phase-aware reactive migration alone. The system degrades gracefully when predictions are unreliable. Future work includes real-system CXL validation, more sophisticated prediction models, and extension to multi-tier hierarchies.

%% ============================================================
%% REFERENCES
%% ============================================================
\bibliographystyle{plainnat}
\bibliography{references}

\end{document}
